% =============================================================================
% Medapati Nisanth Reddy — resume
%
% Compile with pdfLaTeX (Overleaf: Menu -> Compiler -> pdfLaTeX).
%
% ATS notes, so the next edit does not undo them:
%
%  1. \pdfgentounicode=1 writes a ToUnicode map into the PDF. Without it, text
%     extraction from a pdfLaTeX document can return mojibake for perfectly
%     ordinary characters — the single highest-impact ATS fix in this file.
%  2. No icon fonts. FontAwesome glyphs sit in a private-use Unicode range, so a
%     parser reading "\faPhone +91 …" gets an unmapped codepoint glued to the
%     phone number and can drop the field entirely. Contact details are plain
%     text, separated by a real vertical bar.
%  3. Real bullet characters (\textbullet), not label={}. Parsers use bullets to
%     segment responsibilities; unmarked list items merge into one blob.
%  4. One column throughout, no \colorbox behind headings, no tables. Standard
%     section names (EXPERIENCE, EDUCATION, TECHNICAL SKILLS, PROJECTS) because
%     parsers match on those exact words.
%  5. Dates on the right via \hfill, which still extracts in reading order.
%
% Search for VERIFY and TODO before sending this anywhere.
% =============================================================================

\documentclass[11pt,a4paper]{article}

\usepackage[utf8]{inputenc}
\usepackage[T1]{fontenc}
\usepackage[margin=0.55in]{geometry}
\usepackage{xcolor}
\usepackage{enumitem}
\usepackage{parskip}
\usepackage[hidelinks]{hyperref}

% Machine-readable text extraction. Guarded so the file still compiles under
% XeLaTeX or LuaLaTeX, where this primitive does not exist.
\ifdefined\pdfgentounicode
  \input{glyphtounicode}
  \pdfgentounicode=1
\fi

\definecolor{primary}{HTML}{1F1F1F}
\definecolor{secondary}{HTML}{5A6169}
\definecolor{accent}{HTML}{0B5ED7}

\hypersetup{
  colorlinks=true,
  urlcolor=accent,
  linkcolor=primary,
  pdftitle={Medapati Nisanth Reddy - Software Developer},
  pdfauthor={Medapati Nisanth Reddy},
  pdfkeywords={Software Developer, SDLC, Agile, AWS, Docker, Kubernetes, Terraform,
    CI/CD, Python, Java, JavaScript, TypeScript, Kotlin, SQL, React, Next.js, DevOps}
}

\pagestyle{empty}
\setlength{\parindent}{0pt}

% ---------------------------------------------------------------- commands ---

\newcommand{\resumeSection}[1]{%
  \vspace{7pt}%
  {\large\bfseries\color{primary}\MakeUppercase{#1}}\par
  \vspace{1pt}%
  \textcolor{secondary}{\rule{\textwidth}{0.6pt}}\par
  \vspace{3pt}%
}

% #1 title  #2 dates (right)  #3 organisation  #4 detail (right)
%
% One command for every section, so Experience, Education and Projects cannot drift
% into different argument orders — which is exactly what the previous version did:
% arg 2 was the date under EXPERIENCE and the institution under EDUCATION.
%
% \parfillskip=0pt is load-bearing. A paragraph's final line normally ends with
% \parfillskip (0pt plus 1fil), so an \hfill before \par has to share the available
% stretch with it fifty-fifty — the date lands somewhere near the middle of the line
% instead of on the right margin. Zeroing it inside the group gives \hfill all of it.
\newcommand{\entryLine}[4]{%
  \par\vspace{3pt}%
  \begingroup\parfillskip=0pt\relax
    \noindent\textbf{#1}\hfill\textcolor{secondary}{#2}\par
    \noindent\textit{\textcolor{accent}{#3}}\hfill\textit{\textcolor{secondary}{#4}}\par
  \endgroup
  \vspace{1pt}%
}

\newenvironment{bullets}{%
  \begin{itemize}[leftmargin=1.4em, label=\textbullet, itemsep=1.5pt,
                  parsep=0pt, topsep=3pt, partopsep=0pt]
}{%
  \end{itemize}%
  \vspace{2pt}%
}

\newcommand{\skillRow}[2]{\textbf{#1:} #2\par\vspace{2pt}}

\begin{document}

% ------------------------------------------------------------------ header ---
\begin{center}
  {\Huge\bfseries MEDAPATI NISANTH REDDY}\par
  \vspace{3pt}
  {\large Software Developer}\par
  \vspace{5pt}
  % Split over two lines deliberately. Five items on one line wraps somewhere LaTeX
  % chooses, and a contact line that breaks in the middle of a URL looks careless.
  % VERIFY: add your city on the first line if you want one, e.g. "Hyderabad, India \textbar{}"
  +91 8790191677 \textbar{}
  \href{mailto:nisanthreddymedapati@gmail.com}{nisanthreddymedapati@gmail.com} \\[3pt]
  \href{https://nisanth2003.github.io/Portfolio/}{nisanth2003.github.io/Portfolio} \textbar{}
  \href{https://www.linkedin.com/in/nisanth-reddy/}{linkedin.com/in/nisanth-reddy} \textbar{}
  \href{https://github.com/Nisanth2003}{github.com/Nisanth2003}
\end{center}

\vspace{2pt}

% ----------------------------------------------------------------- summary ---
\resumeSection{Professional Summary}

Software developer with over a year of professional experience across the full
\textbf{SDLC} in \textbf{Agile} teams --- requirements gathering, low-level design, code
development, test scenarios, debugging, and deployment. Builds and ships cloud-native
services on \textbf{AWS} and \textbf{Google Cloud} using \textbf{Docker},
\textbf{Kubernetes}, \textbf{Terraform}, and \textbf{CI/CD} pipelines, alongside Android
and web development in \textbf{Kotlin}, \textbf{TypeScript}, and \textbf{Python}.
Delivers reliable, cost-effective solutions with a quality-focused approach, takes
ownership of tasks end to end, and upskills continuously to stay current with cloud and
AI tooling.

% ------------------------------------------------------------ technical skills ---
% Kept near the top: keyword-matching parsers weight an explicit skills block
% heavily, and a recruiter scanning for a stack finds it in two seconds.
\resumeSection{Technical Skills}

\skillRow{Languages}{Python, Java, JavaScript, TypeScript, Kotlin, SQL, HTML, CSS}
\skillRow{Frameworks \& Libraries}{React.js, Next.js, Node.js, Express.js, REST APIs, Tailwind CSS, Jetpack Compose}
\skillRow{Cloud \& DevOps}{AWS (EKS, ALB, IAM), Google Cloud, Docker, Kubernetes, Terraform, CI/CD, GitHub Actions, Git, Linux}
\skillRow{Databases}{PostgreSQL, MySQL, MongoDB, Firebase, Firestore}
\skillRow{AI / ML}{Google Gemini API, Prompt engineering, On-device machine learning}
\skillRow{Practices}{SDLC, Agile development, Requirements gathering, Low-level design, Test scenarios, Debugging, Deployment, Code reviews, Issue resolution and mitigation, Technical documentation, Quality-focused development}

% -------------------------------------------------------------- experience ---
\resumeSection{Experience}

% Dates confirmed 2025: internship from Jun 2025, converted to full-time Dec 2025.
% The portfolio sheet's Experience tab now matches, so the two documents agree.
\entryLine{Software Developer}{Dec 2025 -- Present}{Cloudtechner}{Gurgaon, India (Remote)}
\begin{bullets}
  \item Own delivery of assigned modules through the full \textbf{SDLC} in \textbf{Agile}
        sprints, from requirements gathering and low-level design through code
        development, test scenarios, debugging, and deployment.
  \item Containerise services with \textbf{Docker} and deploy to \textbf{AWS EKS},
        provisioning infrastructure as code with \textbf{Terraform} so environments are
        repeatable and cost-effective to run.
  \item Maintain \textbf{CI/CD} pipelines in \textbf{GitHub Actions}, replacing manual
        deployment steps with automated builds, tests, and releases.
  \item Diagnose and resolve production issues, applying systematic
        \textbf{problem-solving} for fast issue resolution and mitigation.
  \item Write technical documentation for deployment and troubleshooting procedures,
        supporting knowledge sharing across the team.
  % TODO: one bullet with a real number. Anything measurable works — deploy time,
  % build minutes saved, environments provisioned, incidents closed, cost reduced.
  % A single concrete figure outperforms every keyword in this file.
\end{bullets}

\entryLine{DevOps / AI Engineer Intern}{Jun 2025 -- Dec 2025}{Cloudtechner}{Gurgaon, India (Remote)}
\begin{bullets}
  \item Converted to a full-time Software Developer role in December 2025 on the strength
        of internship delivery.
  \item Built and tested application components to specification, taking part in code
        reviews and Agile ceremonies alongside the engineering team.
  \item Automated repetitive deployment and environment-setup tasks, improving day-to-day
        team efficiency and reducing hand-run steps.
  \item Engaged in continuous upskilling across DevOps tooling and cloud services,
        applying each addition to live work rather than leaving it theoretical.
\end{bullets}

% ---------------------------------------------------------------- projects ---
% Three projects, newest first. Named to match the repositories exactly — Hands-Free,
% ClockAlert, Portfolio — so the resume and GitHub line up under a recruiter's eye.
%
% NOTE: the Hands-Free and ClockAlert repositories are currently EMPTY. Push the code
% before you send this out, or both links land on GitHub's "this repository is empty"
% page, which reads worse than having no link at all.
%
% Dropped to keep this to three: EKS AI Automation Pipeline and AI Mock Interview. The
% EKS one is the only cloud/DevOps project you had — if you are applying for anything
% DevOps-leaning, put it back as a fourth entry; the block is in this file's git history.
\resumeSection{Projects}

\entryLine{Hands-Free}{2026}{Kotlin, Android, MediaPipe, CameraX, AccessibilityService, DataStore}{\href{https://github.com/Nisanth2003/Hands-Free}{github.com/Nisanth2003/Hands-Free}}
\begin{bullets}
  \item Android app that controls the phone entirely through hand gestures seen by the
        front camera, using \textbf{MediaPipe} 21-point hand tracking over a
        \textbf{CameraX} pipeline at 20--30 FPS on mid-range devices.
  \item Wrote a geometry-based gesture classifier instead of training a model per gesture,
        keeping the recogniser deterministic, debuggable, and tunable at runtime.
  \item Drives system-wide taps, scrolling, and Back / Home / Recents through an
        \textbf{AccessibilityService}, with a cursor overlay that works over any app
        without requesting \texttt{SYSTEM\_ALERT\_WINDOW}.
  \item Built for accessibility: dwell-to-click for users who cannot coordinate a tap
        gesture, a practice mode, and anti-misfire cooldowns for tremor-prone hands.
\end{bullets}

\entryLine{ClockAlert}{2026}{Kotlin, Android, Room, Geofencing, NotificationListenerService, ML Kit}{\href{https://github.com/Nisanth2003/ClockAlert}{github.com/Nisanth2003/ClockAlert}}
\begin{bullets}
  \item Android alarm app whose alarms fire on a real-world event rather than a clock
        time --- arriving at a destination, or a tracked task finishing.
  \item Designed it around one scheduled alarm per trigger, refined by sparse signals the
        OS pushes in (geofence transitions, notification updates), so there is no
        background polling and no always-on location anywhere in the app.
  \item Guaranteed fallback: if a signal never arrives or a permission is revoked, the
        alarm still rings at the last estimate --- degrading to an ordinary alarm instead
        of silence.
  \item \textbf{Room} persistence with exported schemas and real additive migrations, no
        destructive fallback, across 100+ Kotlin source files.
\end{bullets}

% The live site is in the header, so this slot carries the source — an engineer reading
% the resume wants the repository, and the recruiter already has the link at the top.
\entryLine{Portfolio Site}{2026}{Next.js, TypeScript, Three.js, Google Sheets API, GitHub Actions}{\href{https://github.com/Nisanth2003/Portfolio}{github.com/Nisanth2003/Portfolio}}
\begin{bullets}
  \item Live at \href{https://nisanth2003.github.io/Portfolio/}{nisanth2003.github.io/Portfolio}.
        Statically generated at build time from a private Google Sheet through a service
        account, so publishing new work is a spreadsheet edit rather than a code change or
        a hand-run deploy.
  \item Validates every sheet row at build time and fails the build with the offending row
        number, so malformed data can never reach production.
  \item Deploys to GitHub Pages automatically through \textbf{GitHub Actions}, with a
        committed fallback snapshot so an upstream outage cannot break a release.
\end{bullets}

% --------------------------------------------------------------- education ---
\resumeSection{Education}

% VERIFY the two school years below — inferred from the B.Tech timeline, not confirmed.
\entryLine{B.Tech (Hons.), Computer Science and Engineering}{2021 -- 2025}{Lovely Professional University}{CGPA: 7.8 / 10}
\begin{bullets}
  \item Capstone project: BeWell --- an Android wellness application built with Kotlin and
        Firebase, delivered in Agile sprints.
\end{bullets}

\entryLine{Class XII (Senior Secondary)}{2021}{Sri Chaitanya Bhaskar Bhavan}{90.4\%}
\entryLine{Class X (Secondary)}{2019}{Nalanda Vidya Niketan}{86.4\%}

% -------------------------------------------------------------- additional ---
\resumeSection{Additional Information}

\skillRow{Strengths}{Ownership, Prioritisation, Deadline management, Collaboration, Communication, Proactive issue handling, Upskilling / reskilling / cross-skilling}
\skillRow{Availability}{Flexible to work in shifts; open to relocation and to working from any location}
\skillRow{Academic standing}{No standing arrears; consistent academic performance}

\end{document}
